             11th International Mathematical Competition for University Students
                                   Skopje, 25–26 July 2004

                                      Solutions for problems on Day 2


1. Let A be a real 4 × 2 matrix and B be a real 2 × 4 matrix such that
                                                              
                                              1    0 −1 0
                                           0      1    0 −1
                                     AB = −1 0
                                                               .
                                                        1    0
                                              0 −1 0         1
Find BA.     [20 points]
                    
                     A1               
Solution. Let A =        and B = B1 B2 where A1 , A2 , B1 , B2 are 2 × 2 matrices. Then
                     A2
                                     
                         1   0 −1 0                                          
                       0    1   0 −1    A 1
                                                                A 1 B1  A 1 B2
                       
                       −1 0
                                      =         B1 B2 =
                                 1  0    A2                     A2 B1 A2 B2
                         0 −1 0     1
therefore, A1 B1 = A2 B2 = I2 and A1 B2 = A2 B1 = −I2 . Then B1 = A−1           −1          −1
                                                                     1 , B2 = −A1 and A2 = B2 =
−A1 . Finally,                                                           
                                       A1                              2 0
                         BA = B1 B2           = B1 A1 + B2 A2 = 2I2 =
                                         A2                             0 2


2. Let f, g : [a, b] → [0, ∞) be continuous and non-decreasing functions such that for each x ∈ [a, b] we
have                                    Z xp            Z xp
                                             f (t) dt ≤      g(t) dt
                                            a                    a
    Rbp            Rbp
and a f (t) dt = a g(t) dt.
              Rbp               Rbp
   Prove that a 1 + f (t) dt ≥ a 1 + g(t) dt. [20 points]
                      Rxp                     Rxp
Solution. Let F (x) = a f (t) dt and G(x) = a g(t) dt. The functions F, G are convex, F (a) = 0 =
G(a) and F (b) = G(b) by the hypothesis. We are supposed to show that
                              Z bq                  Z bq
                                             2                    2
                                          0
                                   1 + F (t) dt ≥        1 + G0 (t) dt
                                  a                             a

i.e. The length ot the graph of F is ≥ the length of the graph of G. This is clear since both functions are
convex, their graphs have common ends and the graph of F is below the graph of G — the length of the
graph of F is the least upper bound of the lengths of the graphs of piecewise linear functions whose values
at the points of non-differentiability coincide with the values of F , if a convex polygon P1 is contained in
a polygon P2 then the perimeter of P1 is ≤ the perimeter of P2 .

3. Let D be the closed unit disk in the plane, and let p1 , p2 , . . . , pn be fixed points in D. Show that there
exists a point p in D such that the sum of the distances of p to each of p1 , p2 , . . . , pn is greater than or
equal to 1. [20 points]
Solution. considering as vectors, thoose p to be the unit vector which points into the opposite direction as
Pn
   pi . Then, by the triangle inequality,
i=1

                               n
                               X                         n
                                                         X                n
                                                                          X
                                      |p − pi | ≥ np −         pi = n +         pi ≥ n..
                               i=1                       i=1              i=1
4. For n ≥ 1 let M be an n × n complex matrix with distinct eigenvalues λ1 , λ2 , . . . , λk , with multiplicities
m1 , m2 , . . . , mk , respectively. Consider the linear operator LM defined by LM (X) = M X + XM T , for any
complex n × n matrix X. Find its eigenvalues and their multiplicities. (M T denotes the transpose of M ;
that is, if M = (mk,l ), then M T = (ml,k ).) [20 points]
Solution. We first solve the problem for the special case when the eigenvalues of M are distinct and all sums
λr + λs are different. Let λr and λs be two eigenvalues of M and ~vr , ~vs eigenvectors associated to them, i.e.
                                                                                          T
M~vj = λ~vj for j = r, s. We have M~vr (~vs )T +~vr (~vs )T M T = (M~vr )(~vs )T +~vr M~vs = λr~vr (~vs )T +λs~vr (~vs )T ,
so ~vr (~vs ) is an eigenmatrix of LM with the eigenvalue λr + λs .
     Notice that if λr 6= λs then vectors ~u, w~ are linearly independent and matrices ~u(w)    ~ T and w(~ ~ u)T are
linearly independent, too. This implies that the eigenvalue λr + λs is double if r 6= s.
     The map LM maps n2 –dimensional linear space into itself, so it has at most n2 eigenvalues. We already
found n2 eigenvalues, so there exists no more and the problem is solved for the special case.
     In the general case, matrix M is a limit of matrices M1 , M2 , . . . such that each of them belongs to the
special case above. By the continuity of the eigenvalues we obtain that the eigenvalues of LM are
     • 2λr with multiplicity m2r (r = 1, . . . , k);
     • λr + λs with multiplicity 2mr ms (1 ≤ r < s ≤ k).
(It can happen that the sums λr + λs are not pairwise different; for those multiple values the multiplicities
should be summed up.)


5. Prove that                         Z 1Z 1
                                                            dx dy
                                                                        ≤ 1. [20 points]
                                       0        0   x−1 + | ln y| − 1
Solution 1. First we use the inequality

                                                    x−1 − 1 ≥ | ln x|, x ∈ (0, 1],

which follows from
                                                (x−1 − 1) x=1 = | ln x||x=1 = 0,
                                                  1     1
                                 (x−1 − 1)0 = − 2 ≤ − = | ln x|0 , x ∈ (0, 1].
                                                 x      x
Therefore             Z 1Z 1                     Z 1Z 1                    Z 1Z 1
                                  dx dy                     dx dy                   dx dy
                              −1
                                               ≤                         =                    .
                       0  0 x    + | ln y| − 1    0  0 | ln x| + | ln y|    0  0 | ln(x · y)|

Substituting y = u/x, we obtain
                    Z 1Z 1               Z 1 Z 1            Z 1
                             dx dy               dx    du                      du
                                       =                    =     | ln u| ·         = 1.
                     0  0 | ln(x · y)|    0    u x  | ln u|    0            | ln u|

Solution 2. Substituting s = x−1 − 1 and u = s − ln y,
           Z 1Z 1                     Z ∞Z ∞                  Z ∞ Z u             −u
                       dx dy                    es−u                     es        e
                   −1
                                    =                2
                                                       duds =                 2
                                                                                ds     dsdu.
            0  0 x    + | ln y| − 1    0  s  (s + 1) u         0    0 (s + 1)       u
                          e  s
     Since the function (s+1)2 is convex,

                                           Z u
                                                       es                    eu
                                                                                    
                                                                  u
                                                             ds ≤                  +1
                                            0       (s + 1)2      2       (u + 1)2
so
       Z 1Z 1                         Z ∞                           −u        Z ∞            Z ∞
                                                           eu
                                                                                                      
                     dx dy                  u                       e        1         du           −u
                                  ≤                              +1     du =                 +     e du = 1.
        0   0   x−1 + | ln y| − 1     0     2           (u + 1)2     u       2   0  (u + 1)2    0
6. For n ≥ 0 define matrices An and Bn as follows: A0 = B0 = (1) and for every n > 0
                                                                     
                                   An−1 An−1                 An−1 An−1
                            An =                 and Bn =                 .
                                   An−1 Bn−1                 An−1   0

Denote the sum of all elements of a matrix M by S(M ). Prove that S(Ank−1 ) = S(Akn−1 ) for every n, k ≥ 1.
[20 points]
Solution. The quantity S(Ak−1
                            n ) has a special combinatorical meaning. Consider an n × k table filled with
0’s and 1’s such that no 2 × 2 contains only 1’s. Denote the number of such fillings by Fnk . The filling of
each row of the table corresponds to some integer ranging from 0 to 2n − 1 written in base 2. Fnk equals
to the number of k-tuples of integers such that every two consecutive integers correspond to the filling of
n × 2 table without 2 × 2 squares filled with 1’s.
    Consider binary expansions of integers i and j in in−1 . . . i1 and jn jn−1 . . . j1 . There are two cases:

   1. If in jn = 0 then i and j can be consecutive iff in−1 . . . i1 and jn−1 . . . j1 can be consequtive.

   2. If in = jn = 1 then i and j can be consecutive iff in−1 jn−1 = 0 and in−2 . . . i1 and jn−2 . . . j1 can be
      consecutive.

Hence i and j can be consecutive iff (i + 1, j + 1)-th entry of An is 1. Denoting this entry by ai,j , the sum
          P2n −1         P2n −1
S(Ak−1
   n ) =     i1 =0 · · ·
                                                                                                              k−1
                          ik =0 ai1 i2 ai2 i3 · · · aik−1 ik counts the possible fillings. Therefore Fnk = S(An ).
   The the obvious statement Fnk = Fkn completes the proof.
